% Sections 1--3 of the revised manuscript (r2, six-section structure:
% 1 Introduction, 2 Capability set and sign condition, 3 Market clearing
% formulation, 4 Case study, 5 Results, 6 Conclusion). Numbers that depend on the
% regenerated results are marked \todo{...} and filled from study/results.

\section{Introduction}
\label{sec:intro}

As synchronous generation is displaced, inertia and fast reserve have to be
procured explicitly, and increasingly from converter-interfaced
resources~\cite{ulbig2014inertia,tielens2016inertia,milano2018lowinertia}.
Grid-forming (GFM) control, in which the converter behaves as a voltage source
behind an impedance, is the main technical
response~\cite{zhong2011synchronverter,rocabert2012control,dorfler2023control},
and system operators now specify GFM capability rather than encourage
it~\cite{matevosyan2019gfm,lin2020roadmap,entsoe2020gfm}. The energy a GFM
converter sells, the reserve it holds and the inertial power it injects after
a disturbance all pass through the same semiconductor bridge, so they are
limited jointly, and the limit is a current~\cite{yazdani2010vsc}.

Scheduling formulations have treated the frequency side of this problem in
increasing detail. Governor rate limits entered the optimal power flow
in~\cite{chavez2014governor}; inertia became a decision variable
in~\cite{teng2016stochastic}; frequency services were
co-optimised in~\cite{badesa2019simultaneous,trovato2019unit}; and
requirement uncertainty and the siting of virtual inertia were addressed
in~\cite{paturet2020stochastic,poolla2019placement}. On the supply side these
models bound what a resource can offer by active power. Three recent
formulations require converter inertia to be backed by unused
headroom~\cite{park2026inertia,jiang2025frequency,beigifard2026frequency},
and a fourth embeds a learned frequency surrogate in the day-ahead
schedule~\cite{jiang2026learning}. They differ in most respects, but in each
the converter capability has the form
%
\begin{equation}
  P_i + R_i + H_i \;\le\; \bar{P}_i ,
  \label{eq:box}
\end{equation}
%
with $P_i$ the energy schedule, $R_i$ the reserve and $H_i$ the inertial
power. None contains a current limit, a reactive-power term or the terminal
voltage (Table~\ref{tab:related}).

System operators describe the capability differently. AEMO specifies overload
capability as a level with a duration, for example \SI{150}{\percent} of rated
current for \SI{2}{\second}, and reports that the active headroom of an
inverter depends on its pre-disturbance operating
point~\cite{aemo2023gfm,aemo2024inertia}; MISO's draft specification addresses
operation at short-duration current limits~\cite{miso2024gfm}. At device and
feeder level the interaction between reactive support and active reserve
through the apparent-power limit is
established~\cite{energies2020reserve}, and current-based inverter limits
have been used in photovoltaic--storage scheduling and in batteries providing
frequency and voltage control
together~\cite{polat2025inverter,zecchino2021concurrent}. The remaining gap is
a market clearing in which the current limit is shared between energy, reserve
and inertia, the requirement multipliers are prices, and the effect of the
representation is quantified.

\begin{table*}[t]
  \caption{Scheduling formulations with shared converter capability and what
  each represents ($\yes$ present, $\no$ absent, --- not applicable). Entries
  follow the constraint sets stated in each paper.}
  \label{tab:related}
  \centering
  \footnotesize
  \begin{tabular*}{\textwidth}{@{\extracolsep{\fill}}l l c c c c c c c l@{}}
    \toprule
    & & Shared & Current & Reactive & Terminal & Short-term
    & Storage & Prices & Frequency \\
    Reference & Test system & capability & limit & power & voltage
    & $\neq$ continuous & state & as duals & model \\
    \midrule
    Park et al.~\cite{park2026inertia} & IEEE 30
      & \yes & \no & \no & \no & \no & \yes & \yes & --- \\
    Jiang and Li~\cite{jiang2025frequency} & WSCC 9, IEEE 39
      & \yes & \no & \no & \no & \no & \no & \no & EMT surrogate \\
    Jiang and Li~\cite{jiang2026learning} & IEEE 9
      & \yes & \no & \no & \no & \no & \yes & \no & EMT surrogate \\
    Beigi Fard et al.~\cite{beigifard2026frequency} & IEEE 9, 118
      & \yes & \no & \no & \no & \no & \yes & \no & PWL nadir \\
    Garozzo and Tina~\cite{energies2020reserve} & inverter, LV feeder
      & --- & \yes & \yes & \no & \no & \yes & --- & --- \\
    \addlinespace
    This work & IEEE 39, RTS-24
      & \yes & \yes & \yes & \yes & \yes & \yes & \yes & RMS check \\
    \bottomrule
  \end{tabular*}
\end{table*}

The physical limit on a converter is
%
\begin{equation}
  \sqrt{P^2 + Q^2} \;\le\; V I_{\max} S ,
  \label{eq:disc}
\end{equation}
%
a disc in the $P$--$Q$ plane whose radius scales with the terminal voltage
$V$. Equation~\eqref{eq:box} is the projection of Eq.~\eqref{eq:disc} onto
$Q = 0$ at $V = \SI{1.0}{\pu}$ with $I_{\max}$ fixed at the continuous rating.
It therefore omits the short-term current rating, which understates the
capability available for an inertial response; omits reactive current, which
overstates capability whenever $Q \neq 0$; and fixes the voltage, which
overstates capability below \SI{1.0}{\pu} and understates it above. The three
errors do not have the same sign, so a constant derating of
Eq.~\eqref{eq:box} cannot remove them together.

The contributions of this paper are as follows.
\begin{enumerate}
  \item The continuous and short-term current ratings, reactive current and
        terminal voltage of a GFM converter are represented as a polyhedral
        capability set and integrated in a multi-period joint clearing of
        energy, reserve and inertia with marginal pricing. The clearing
        remains a linear programme, so the prices are dual variables.
  \item A piecewise condition is derived for the sign of the error made by
        the active-power bound of Eq.~\eqref{eq:box}. The condition depends
        on the short-term radius, the bound and the reactive current held
        during the response, and not on the sustained active power.
  \item The formulation is evaluated on the IEEE 39-bus and RTS-24 systems
        against three comparison models and an exact second-order cone
        benchmark, under alternative converter rating conventions and
        randomised converter placements. The final schedules are validated
        by ac power flow against voltage, generator reactive, branch and
        balance limits; the procured response is constrained to be
        deliverable after the largest generator outages; and the scheduled
        inertial power is checked against a current-limited RMS converter
        model.
\end{enumerate}

Fig.~\ref{fig:overview} summarises the approach. The formulation extends the
shared-capability clearing of~\cite{park2026inertia}, whose treatment of
inertia as a priced power service and of storage energy recovery is retained;
what changes is the shape of the set the services share.

\begin{figure*}[t]
  \centering
  \includegraphics[width=\textwidth]{fig1_overview}
  \caption{Overview: physical current limit, polyhedral capability set,
  multi-period clearing with network and ac validation, and procurement
  outcomes.}
  \label{fig:overview}
\end{figure*}

\section{Capability set and sign condition}
\label{sec:capability}

\subsection{Continuous and short-term limits}

Converter $i$ has apparent-power rating $S_i$, terminal voltage $V_i$ (pu),
continuous current limit $I^{\mathrm{c}}$ and short-term current limit
$I^{\mathrm{s}} \ge I^{\mathrm{c}}$ (pu of rated current). Reserve that is
called must be sustained for the containment period, whereas the inertial
response lasts of the order of a second. Two constraints follow:
%
\begin{align}
  \sqrt{(P_i + R_i)^2 + Q_i^2} &\;\le\; V_i I^{\mathrm{c}} S_i ,
  \label{eq:cont} \\
  \sqrt{(P_i + \chi R_i + H_i)^2 + \tilde{Q}_i^2}
    &\;\le\; V^{\mathrm{sh}}_i \big[ I^{\mathrm{c}}
      + a_i (I^{\mathrm{s}} - I^{\mathrm{c}}) \big] S_i .
  \label{eq:short}
\end{align}
%
Equation~\eqref{eq:cont} applies the continuous (thermal) rating to the
sustained schedule. It is imposed at both ends of the reserve activation,
$(P_i, Q_i)$ and $(P_i + R_i, Q_i)$, because a charging converter
($P_i < 0$) carries its largest current before its reserve is called; the set
is convex, so the intermediate points are covered.
Equation~\eqref{eq:short} applies the short-term rating at the instant of the
inertial response. The factor $\chi \in [0,1]$ is the fraction of the reserve
already activated at that instant and is set to one, the simultaneous
envelope. $V^{\mathrm{sh}}_i$ is the terminal voltage during the event; it
equals $V_i$ unless the correction of Section~\ref{sec:applied} is applied.
The band availability $a_i \in [0,1]$ is defined below.
Fig.~\ref{fig:capability} compares the two constraints with
Eq.~\eqref{eq:box}.

\begin{figure}[t]
  \centering
  \includegraphics[width=\columnwidth]{fig2_capability}
  \caption{Capability set and active-power bound at
  $I^{\mathrm{c}} = \SI{1.0}{\pu}$, $I^{\mathrm{s}} = \SI{1.5}{\pu}$.}
  \label{fig:capability}
\end{figure}

\subsection{Reactive current priority}

During a short-term excursion the current-priority setting of the converter
determines how much reactive current is
retained~\cite{qoria2020current,bottrell2014comparison}. The setting is a
property of the converter, so it enters as a parameter:
%
\begin{equation}
  \tilde{Q}_i = \alpha_i Q_i , \qquad \alpha_i \in [0,1] ,
  \label{eq:priority}
\end{equation}
%
with $\alpha_i = 1$ for reactive priority and $\alpha_i = 0$ for ideal active
priority. Writing the relation as a bound $0 \le \tilde{Q}_i \le Q_i$ would
let the optimisation choose the converter's current allocation.

\subsection{Storage state of charge}

A store limits the services in two ways. The short-term band is available
only above an energy floor $\underline{\sigma}^{\mathrm{e}}$ and fully
available above a breakpoint $\sigma^{\star}$,
%
\begin{equation}
  0 \;\le\; a_i^t \;\le\; \min\!\left(1, \;
    \frac{\min\!\left(\sigma_i^{t-1}, \sigma_i^{t}\right)
      - \underline{\sigma}^{\mathrm{e}}}
    {\sigma^{\star} - \underline{\sigma}^{\mathrm{e}}}\right) ,
  \label{eq:gate}
\end{equation}
%
and the energy behind the services must be held for their delivery windows
$\tau^{\mathrm{R}}$ and $\tau^{\mathrm{H}}$,
%
\begin{equation}
  \frac{\tau^{\mathrm{R}} R_i^t + \tau^{\mathrm{H}} H_i^t}{\eta_{\mathrm{d}}}
    \;\le\; \Bigl( \min\!\left(\sigma_i^{t-1}, \sigma_i^{t}\right)
      - \underline{\sigma}^{\mathrm{e}} \Bigr) E_i .
  \label{eq:deliverability}
\end{equation}
%
Both are evaluated at the two ends of the period and are linear in the state
of charge. The linear ramp in Eq.~\eqref{eq:gate} is an assumption; moving the
breakpoint $\sigma^{\star}$ across its range changes the reported cost
differences by less than \num{0.05} percentage points.

\subsection{Polyhedral form}
\label{sec:polygon-error}

Equations~\eqref{eq:cont} and~\eqref{eq:short} are second-order cones. Each
is replaced by an inscribed regular $m$-gon~\cite{bental2001polyhedral},
%
\begin{align}
  (P_i + R_i)\cos\theta_k + Q_i \sin\theta_k
    &\;\le\; V_i I^{\mathrm{c}} S_i \cos\tfrac{\pi}{m} ,
  \label{eq:polygon} \\
  (P_i + \chi R_i + H_i)\cos\theta_k + \tilde{Q}_i \sin\theta_k
    - a_i \beta_i
    &\;\le\; V^{\mathrm{sh}}_i I^{\mathrm{c}} S_i \cos\tfrac{\pi}{m} ,
  \label{eq:polygon-short}
\end{align}
%
for $\theta_k = 2\pi k / m$, $k = 0, \dots, m-1$, with
$\beta_i = V^{\mathrm{sh}}_i (I^{\mathrm{s}} - I^{\mathrm{c}}) S_i
\cos(\pi/m)$ the width of the overload band. The band width is a parameter,
so both constraints are affine. The polygon is an inner approximation: the
clearing cost is an upper bound on that of the exact conic problem. With
$m = 24$ the cleared cost exceeds the exact second-order cone solution by at
most \SI{0.066}{\percent} over the configurations of
Section~\ref{sec:results}, and by \SI{0.0045}{\percent} with $m = 96$. The case
study uses $m = 24$, and $m = 192$ where the polygon error would be comparable
to the effect measured.

\subsection{Sign of the error of the active-power bound}
\label{sec:sign}

Consider a sustained operating point admitted by both representations. The
radius allowed by the short-term constraint is
%
\begin{equation}
  L_i = V^{\mathrm{sh}}_i \big[ I^{\mathrm{c}}
        + a_i (I^{\mathrm{s}} - I^{\mathrm{c}}) \big] S_i .
  \label{eq:radius}
\end{equation}
%
The active-power bound leaves inertial headroom
$H^{\mathrm{ap}}_i = \bar{P}_i - (P_i + R_i)$ and the capability set leaves
$H^{\mathrm{set}}_i = \sqrt{L_i^2 - \tilde{Q}_i^2} - (P_i + \chi R_i)$, so
%
\begin{equation}
  \Delta H_i = H^{\mathrm{set}}_i - H^{\mathrm{ap}}_i
    = \sqrt{L_i^2 - \tilde{Q}_i^2} - \bar{P}_i + (1 - \chi) R_i .
  \label{eq:difference}
\end{equation}
%
At $\chi = 1$ the difference does not depend on the sustained active power.

\medskip
\noindent\textbf{Proposition.} Let $\chi = 1$ and $L_i$ be given by
Eq.~\eqref{eq:radius}.
\emph{(i)} If $L_i > \bar{P}_i$, define
$\tilde{Q}^{\mathrm{crit}}_i = \sqrt{L_i^2 - \bar{P}_i^2}$. The active-power
bound understates the inertial headroom for
$|\tilde{Q}_i| < \tilde{Q}^{\mathrm{crit}}_i$, is exact at
$|\tilde{Q}_i| = \tilde{Q}^{\mathrm{crit}}_i$, and overstates it for
$|\tilde{Q}_i| > \tilde{Q}^{\mathrm{crit}}_i$.
\emph{(ii)} If $L_i = \bar{P}_i$, the two agree only at $\tilde{Q}_i = 0$ and
the bound overstates elsewhere.
\emph{(iii)} If $L_i < \bar{P}_i$, the bound overstates for every admissible
$\tilde{Q}_i$.
\medskip

The proof follows from Eq.~\eqref{eq:difference}. A higher $I^{\mathrm{s}}$, a
higher voltage or a larger band availability increases $L_i$ and moves the
bound towards understatement; a larger $|\tilde{Q}_i|$ moves it towards
overstatement. The multiplier that would make Eq.~\eqref{eq:box} exact,
$\gamma_i = \sqrt{L_i^2 - \tilde{Q}_i^2} / \bar{P}_i$, varies with voltage,
reactive dispatch, short-term rating and band availability, so no constant
derating reproduces it. The quantities in the condition are available at the
time of clearing: terminal voltage from the state estimator, ratings from
equipment data and reactive dispatch from the schedule. Region (iii) applies
when $V^{\mathrm{sh}}_i I^{\mathrm{s}} < \bar{P}_i / S_i$ with the band
withdrawn, i.e.\ at low voltage with little overload capability.

Fig.~\ref{fig:signmap} gives the crossing for $\bar{P} = S$ and full
band availability. At $I^{\mathrm{s}} =
\SI{1.5}{\pu}$ the crossing exceeds the rating over the normal voltage band,
so the bound is conservative for practical reactive dispatch. At
$I^{\mathrm{s}} = \SI{1.1}{\pu}$ it lies between \num{0.30} and \num{0.58}, a
range that ordinary reactive duty reaches. The condition was checked against
the exact disc at \num{1884} points over $V \in [0.90, 1.08]$,
$I^{\mathrm{s}} \in [1.0, 1.8]$, three band availabilities and two bounds,
and gives the correct sign at all of them.

\begin{figure}[t]
  \centering
  \includegraphics[width=\columnwidth]{fig3_sign_map}
  \caption{Crossing $\tilde{Q}^{\mathrm{crit}}$ (pu of rating) at
  $\bar{P} = S$ and full band availability. Below the heavy contour
  $L < \bar{P}$ and the bound overstates at every reactive dispatch.}
  \label{fig:signmap}
\end{figure}

\section{Market clearing formulation}
\label{sec:clearing}

\subsection{Objective and requirements}

The clearing covers $T$ periods of length $\Delta$ with converters $i \in
\mathcal{I}$, synchronous units $g \in \mathcal{G}$ and zones $z \in
\mathcal{Z}$. The decisions are converter discharge $p^{\mathrm{d}}$, charge
$p^{\mathrm{c}}$, curtailment $c$, reactive output $Q$, reserve $R$, inertial
power $H$, band availability $a$ and state of charge $\sigma$, and
synchronous energy $p^{\mathrm{g}}$ and reserve $r^{\mathrm{g}}$. With
$e_i^t$ the available converter injection, the net converter injection is
%
\begin{equation}
  P_i^t = e_i^t + p^{\mathrm{d},t}_i - p^{\mathrm{c},t}_i - c_i^t ,
  \label{eq:injection}
\end{equation}
%
and the objective is
%
\begin{align}
  \min \; \Delta \sum_t \Big[ & \sum_i \big( \kappa^{\mathrm{e}}
      p^{\mathrm{d},t}_i + \kappa^{\mathrm{r}} R_i^t
      + \kappa^{\mathrm{h}} H_i^t + \kappa^{\mathrm{q}} Q_i^t
      + \kappa^{\mathrm{cyc}} ( p^{\mathrm{d},t}_i + p^{\mathrm{c},t}_i )
      \big) \nonumber \\
  & + \sum_g \big( \kappa^{\mathrm{ge}} p^{\mathrm{g},t}_g
    + \kappa^{\mathrm{gr}} r^{\mathrm{g},t}_g \big) \Big] ,
  \label{eq:objective}
\end{align}
%
with energy offers in \$/MWh and availability offers in \$/MW/h. The power
balance and the reserve, inertia and zonal reactive requirements are
%
\begin{align}
  \sum_i P_i^t + \sum_g p^{\mathrm{g},t}_g
    &= D^t + \ell^t(\mathbf{x}^t) ,
  \label{eq:balance} \\
  \sum_i R_i^t + \sum_g r_g^{\mathrm{g},t} &\;\ge\; \rho \, D^t ,
  \label{eq:reserve} \\
  \sum_{i} H_i^t &\;\ge\; P^{\mathrm{loss}}_{\omega}
    - \tfrac{2\dot{f}^{\max}}{f_0}
      \textstyle\sum_{g \ne \omega} \mathcal{H}_g S_g
    \quad \forall\, \omega \in \mathcal{G} , \label{eq:inertia} \\
  \sum_{i \ne \omega} H_i^t &\;\ge\; P^{\mathrm{loss}}_{\omega}
    - \tfrac{2\dot{f}^{\max}}{f_0}
      \textstyle\sum_{g} \mathcal{H}_g S_g
    \quad \forall\, \omega \in \mathcal{I} , \label{eq:inertia-conv} \\
  \sum_{i \in \mathcal{I}_z} Q_i^t &\;\ge\; \varphi \, Q^{z,t}
    \quad \forall\, z \in \mathcal{Z}_{\mathcal{I}} .
  \label{eq:reactive}
\end{align}
%
In Eq.~\eqref{eq:balance}, $\ell^t$ is the linearised loss term of
Section~\ref{sec:acloop}; it is zero in the first (lossless) clearing. The
inertia requirement is written once per credible outage $\omega$, and the
contribution of the tripped resource is removed from the side on which it
appears. The loss $P^{\mathrm{loss}}_\omega$ is the scheduled active bound of
the unit for both fleets, so the clearing cannot reduce its own contingency
by de-loading the unit that sets it. The reactive requirement is written
for the zones $\mathcal{Z}_{\mathcal{I}} \subseteq \mathcal{Z}$ that
contain a converter; the reactive load of the other zones is served by plant
outside the clearing. Commitment is fixed.

The store follows
%
\begin{equation}
  \sigma_i^{t} = \sigma_i^{t-1}
    - \frac{\Delta \, p^{\mathrm{d},t}_i}{\eta_{\mathrm{d}} E_i}
    + \frac{\Delta \, \eta_{\mathrm{c}} \, p^{\mathrm{c},t}_i}{E_i} ,
  \qquad \sigma_i^{T} = \sigma_i^{0} ,
  \label{eq:store}
\end{equation}
%
without a complementarity constraint; simultaneous charge and discharge of a
unit does not occur in any reported schedule. Branch flows are limited through
the shift-factor matrix $\boldsymbol{\Psi}$,
%
\begin{equation}
  \big| \boldsymbol{\Psi} \big( \mathbf{P}^t + \mathbf{p}^{\mathrm{g},t}
  - \mathbf{d}^t \big) \big| \;\le\; \mathbf{f}^{\max} ,
  \label{eq:flows}
\end{equation}
%
and the net converter injection is bounded by the scheduled active limit,
$-P^{\mathrm{sch}}_i \le P_i^t \le P^{\mathrm{sch}}_i$.

Table~\ref{tab:bounds} distinguishes the three bounds that appear in the
comparison. $P^{\mathrm{sch}}_i$ and $S_i$ are common to all models;
$\bar{P}_i$ belongs to the comparison model of Eq.~\eqref{eq:box} only.
Results are reported for both $\bar{P}_i = S_i$ and
$\bar{P}_i = P^{\mathrm{sch}}_i$.

\begin{table}[t]
  \caption{Converter bounds used in the clearing and in the comparison
  model.}
  \label{tab:bounds}
  \centering
  \small
  \begin{tabularx}{\columnwidth}{@{}l X X@{}}
    \toprule
    Bound & Quantity limited & Value \\
    \midrule
    $P^{\mathrm{sch}}_i$ & net scheduled active injection
      & rated active power of the converter \\
    $S_i$ & radius of the capability disc
      & $P^{\mathrm{sch}}_i / (P/S)$, Section~\ref{sec:systems} \\
    $\bar{P}_i$ & $P_i + R_i + H_i$ in Eq.~\eqref{eq:box}
      & $S_i$ or $P^{\mathrm{sch}}_i$ (both reported) \\
    \bottomrule
  \end{tabularx}
\end{table}

\subsection{Prices}
\label{sec:prices}

Prices are the multipliers of the requirements~\cite{schweppe1988spot}:
$\lambda^t$ for Eq.~\eqref{eq:balance}, $\mu^{\mathrm{r},t}$ for
Eq.~\eqref{eq:reserve} and $\mu^{z,t}$ for Eq.~\eqref{eq:reactive}. The
inertia requirement has one row per outage, and the individual rows are
degenerate, so the multiplier of a single row depends on the basis returned by
the solver. The inertia price reported here is therefore defined as the
marginal cost of raising all inertia rows of a period by one megawatt, i.e.\
the sum of their multipliers, averaged over the day. This quantity is a
directional derivative of the optimal cost. It is verified by re-clearing with
perturbations of \SIlist{1;0.25;0.01}{\MW}, and its left and right limits are
reported as a price range. With fixed commitment
the clearing is a linear programme, so no pricing rule for integer decisions
is needed~\cite{oneill2005efficient,gribik2007market}.

\subsection{AC-feasible clearing}
\label{sec:acloop}

Equations~\eqref{eq:flows} and~\eqref{eq:reactive} do not guarantee that a
schedule satisfies the ac network limits. The final schedules are therefore
obtained by successive linear programming~\cite{alsac1990further,
castillo2016successive}. Round $k$ consists of four steps.
\begin{enumerate}
  \item \emph{Clearing.} The linear programme is solved with the rows of step
        4 from round $k-1$. A lexicographic second solve selects, among
        schedules of equal cost, the one with least storage throughput, which
        makes the schedule unique; prices are taken from the first solve.
  \item \emph{Power flow and voltage control.} An ac power flow is solved for
        each period with converters as $PQ$ injections. Generator and
        reference-machine voltage setpoints and transformer taps are adjusted
        by a least-change search until bus voltages and machine reactive
        outputs are within limits; reactive limits are not enforced by the
        solver, so a violation appears as a reactive excess. These controls
        are not market products and carry no price.
  \item \emph{Validation.} The schedule is checked against every bus voltage
        limit, every machine reactive limit, the loading of every branch at
        its more heavily loaded end, the difference between the cleared and
        the actual output of the reference machine together with its active
        limits and sold reserve, and, for the capability set, each converter
        current recomputed at the ac terminal voltage.
  \item \emph{Linearisation.} Sensitivities of bus voltages, machine reactive
        outputs, branch flows and losses to the controlled injections are
        taken about the power-flow solution by central differences. They
        give (a) the loss term $\ell^t$ of Eq.~\eqref{eq:balance}; (b)
        branch apparent-power limits as inscribed polygons, tightened by the
        ratio of measured to linearised loading; (c) rows that remove small
        residual voltage and reactive violations or otherwise prevent them
        from increasing; and (d) a trust region on the active and reactive
        injection at each controlled bus, which shrinks geometrically from
        \SI{4}{\percent} of peak demand. The converter terminal voltages in
        Eqs.~\eqref{eq:polygon}--\eqref{eq:polygon-short} are updated.
\end{enumerate}
The iteration stops when the validation of step 3 passes within the
tolerances of Table~\ref{tab:acfinal}. Every subproblem is a linear
programme, and the reported prices are the multipliers of the final one, i.e.\
marginal prices conditional on the final linearisation and control settings.

\subsection{Post-contingency deliverability}
\label{sec:secure-form}

Equation~\eqref{eq:flows} limits the pre-contingency dispatch only. For a set
$\Omega$ of secured generator outages, the procured response is also required
to reach the network. After the loss of unit $\omega$, its scheduled output
$p^{\mathrm{g},t}_\omega$ is replaced by the held response, deployed in
proportion $\phi^t_\omega \in [0,1]$ to the holdings $R_i^t + H_i^t$ and
$r^{\mathrm{g},t}_g$ ($g \ne \omega$), and the remainder is supplied by the
surviving synchronous machines in proportion to their stored energy
$\mathcal{H}_g S_g$. With $\mathbf{s}_\omega$ that sharing vector and
$\mathbf{u}^t_\omega$ the vector of deployed holdings,
%
\begin{equation}
  \big| \boldsymbol{\Psi} \big( \mathbf{P}^t + \mathbf{p}^{\mathrm{g},t}
    - \mathbf{d}^t
    - p^{\mathrm{g},t}_\omega \mathbf{1}_\omega
    + \phi^t_\omega \mathbf{u}^t_\omega
    + ( p^{\mathrm{g},t}_\omega - \phi^t_\omega
        \mathbf{1}^{\!\top} \mathbf{u}^t_\omega ) \mathbf{s}_\omega
    \big) \big| \;\le\; \varepsilon \, \mathbf{f}^{\max}
  \label{eq:secure}
\end{equation}
%
for all $\omega \in \Omega$ and $t$, where $\varepsilon = 1.25$ is the
short-term emergency rating factor. The fraction
$\phi^t_\omega = \min(1, p^{\mathrm{g},t}_\omega / \mathbf{1}^{\!\top}
\mathbf{u}^t_\omega)$ is a ratio of decisions; it is fixed from the previous
solution and iterated with damping until it changes by less than
\num{e-3}, so Eq.~\eqref{eq:secure} is linear in each solve.
